\subsection{Types, precision, and compiler support}

\verb@__int128@ (also \verb@__int128_t@) is a signed integer in
$[-2^{127},2^{127}-1]$; \verb@unsigned __int128@ (\verb@__uint128_t@)
is in $[0,2^{128}-1]$. Signed overflow is undefined, even with 128 bits.
Widen \emph{before} arithmetic: \verb@__int128(a)*b@, not
\verb@__int128(a*b)@. For differences, widen before subtraction too.

Typical x86-64 GNU C++ floating-point precision:
\begin{center}
\begin{tabular}{l|r|r}
Type & Significand bits & Decimal digits (approx.) \\ \hline
\verb@double@ & 53 & 16 \\
\verb@long double@ & 64 & 19 \\
\verb@__float128@ & 113 & 34
\end{tabular}
\end{center}

\verb@long double@ is platform-dependent: 16 bytes of storage can mean
80-bit extended precision plus padding, not quadruple precision. Check
\verb@numeric_limits<long double>::digits@, or \verb@LDBL_MANT_DIG@
from \verb@<cfloat>@; \verb@FLT128_MANT_DIG@ from \verb@<quadmath.h>@
is 113. \verb@using LD = __float128;@ makes \verb@LD@ quadruple
precision regardless of the alias's name.

Both 128-bit types are compiler extensions, not C++17 additions. They work
with \verb@g++ -std=c++17@ on supported targets; a newer language standard
does not change \verb@long double@ precision. C++23 separately introduced
the optional \verb@std::float128_t@ in \verb@<stdfloat>@. Availability
depends on the compiler, standard library, and target.

\subsection{Input and output}

For the ordinary types, use matching formats:
\begin{center}
\begin{tabular}{l|l|l}
Type & \verb@scanf@ & \verb@printf@ \\ \hline
\verb@long long@ & \verb@%lld@ & \verb@%lld@ \\
\verb@unsigned long long@ & \verb@%llu@ & \verb@%llu@ \\
\verb@double@ & \verb@%lf@ & \verb@%f@ \\
\verb@long double@ & \verb@%Lf@ & \verb@%Lf@
\end{tabular}
\end{center}

These ordinary types support \verb@cin@ and \verb@cout@ directly.
In GNU C++17, the 128-bit extension types have no direct standard stream
overloads. Never pass \verb@__int128@ to \verb@%lld@ or
\verb@__float128@ to \verb@%Lf@: mismatched variadic arguments are invalid.

\subsubsection{Integer conversion}

There is no standard \verb@scanf@ or \verb@printf@ format for
\verb@__int128@. If input values fit in \verb@long long@, read that type
and then widen. Otherwise read a string and parse it. The following helpers
handle the full signed range, including the minimum value. \verb@parse128@
accepts an optional sign and decimal digits; on invalid input or overflow
it returns false without changing the output.

\kactlimport{Int128IO.h}

\begin{lstlisting}[language=C++]
string token;
__int128_t n;
if (cin >> token && parse128(token, n))
    cout << str128(n) << '\n';
\end{lstlisting}

\subsubsection{Quadruple-precision conversion}

Use \verb@strtoflt128@ for input and \verb@quadmath_snprintf@ for output.
Include \verb@<quadmath.h>@ and link with \verb@-lquadmath@:
\begin{lstlisting}[language=bash]
g++ -std=c++17 -O2 solution.cpp -lquadmath
\end{lstlisting}

This complete program assumes a valid floating-point input token:
\begin{lstlisting}[language=C++,label=quadmath-example]
#include <quadmath.h>
#include <iostream>
#include <string>
using namespace std;

int main() {
    string token;
    if (!(cin >> token)) return 0;
    __float128 x = strtoflt128(token.c_str(), nullptr);
    char buf[128];
    int n = quadmath_snprintf(buf, sizeof(buf), "%.36Qg", x);
    if (n < 0 || n >= int(sizeof(buf))) return 1;
    cout << buf << '\n';
}
\end{lstlisting}

\verb@%.36Qg@ prints 36 significant digits, sufficient for binary128
round trips. \verb@%.15Qf@ prints 15 digits after the decimal point;
\verb@%Qe@ uses scientific notation. The format must contain exactly one
conversion, with no surrounding text (not even a newline). Print text and
newlines separately. A fixed-point result can require thousands of digits;
check the returned length and increase the buffer if needed.

Ordinary \verb@printf("%Qf", x)@ only works with some C-library hooks;
do not rely on it. There is no portable matching \verb@scanf@ conversion.
For unchecked input, use \verb@strtoflt128@'s end pointer to check that the
whole token was read, and inspect \verb@errno@ for range errors.

\subsection{Avoiding accidental precision loss}

\begin{itemize}
\item \verb@__float128 x = 0.1;@ first rounds the literal to \verb@double@.
  Prefer \verb@__float128(1)/10@ or \verb@strtoflt128("0.1", nullptr)@.
  GCC's \verb@Q@ literal suffix may need \verb@-fext-numeric-literals@
  in strict C++ mode; the examples above do not need it.
\item Basic 128-bit arithmetic does not require \verb@-lquadmath@.
  Quadmath conversions and functions such as \verb@sqrtq@ and
  \verb@fabsq@ do. Ordinary \verb@cmath@ overloads are not a substitute.
\item \verb@cout << (long double)x@ is an approximate output shortcut,
  not full-precision binary128 output. Cast only at the end, and only if
  the problem's tolerance permits the loss. More printed digits cannot
  recover precision already lost in a conversion or calculation.
\item For integer-coordinate geometry, keep predicates in sufficiently
  wide integer arithmetic and convert only when a division is required.
  Approximate intersection coordinates can lose accuracy again when used
  to calculate a thin triangle's area. More precision helps cancellation;
  it does not make floating-point predicates exact. An algebraic area
  formula may avoid constructing rounded coordinates altogether.
\end{itemize}

Sources: GCC's manuals for
\href{https://gcc.gnu.org/onlinedocs/gcc/_005f_005fint128.html}{128-bit integers},
\href{https://gcc.gnu.org/onlinedocs/gcc/Floating-Types.html}{floating types},
\href{https://gcc.gnu.org/onlinedocs/libquadmath/I_002fO-Library-Routines.html}{quadmath I/O},
and the \href{https://en.cppreference.com/w/cpp/types/floating-point.html}{C++23 fixed-width floating types}.
